\subsection{Measures on the dual of a nuclear space}

Let F be a locally convex space. Let \(T_{s}\) be the weak topology \(\sigma(\mathbf{F}^{\prime},\mathbf{F})\) on \(F^{\prime}\) , and \(T_{c}\) the topology of uniform convergence on the compact convex subsets of F. By Mackey's theorem (TVS, IV, §1, No.~1, Th. 1) the topologies \(T_{s}\) and \(T_{c}\) on \(F^{\prime}\) are compatible with the duality between F and \(F^{\prime}\) ; the same is therefore true of every locally convex topology T on \(F^{\prime}\) intermediate to \(T_{s}\) and \(T_{c}\) . If T is such a topology, and \(F^{\prime}_{T}\) denotes the space \(F^{\prime}\) equipped with T, we shall identify F with the dual of \(F^{\prime}_{T}\) . The promeasures on \(F^{\prime}\) are therefore the same for all topologies T of the preceding type, and if \(\mu\) is such a promeasure then its Fourier transform is a function on F.

One calls Sazonov's topology on F the locally convex topology S defined by the continuous seminorms N satisfying the following condition: \(N^{2}\) is a positive quadratic form on F and there exists a continuous positive quadratic form H on F such that \(\operatorname{Tr}\left(\mathrm{N}^{2}/\mathrm{H}\right)<+\infty\) . The topology S is coarser than the given topology on F; if these topologies are identical, F is said to be nuclear. This class of spaces will be studied in detail later on.

THEOREM 2 (Minlos). --- Let F be a locally convex space, T a locally convex topology on \(F'\) intermediate to \(T_{s}\) and \(T_{c}\) , and \(\mu\) a promeasure on \(F'_{T}\) . Assume that the Fourier transform \(\Phi\) of \(\mu\) is continuous on F for the Sazonov topology. Then \(\mu\) is a measure on \(F'_{T}\) .

Let \(\varepsilon > 0\) . Since \(\Phi\) is continuous for the Sazonov topology on \(F\) , there exist two continuous positive quadratic forms \(Q\) and \(H\) on \(F\) such that \(\operatorname{Tr}(Q / H) < +\infty\) and

\[
\Phi (0) - \mathcal {R} \Phi (x) \leqslant \varepsilon / 6
\]

for every \(x \in \mathbf{F}\) such that \(\mathbf{Q}(x) \leqslant 1\) . By Prop. 8 of No.~8, \(|\mathcal{R}\Phi(x)| \leqslant \Phi(0)\) for all \(x \in \mathbf{F}\) , whence

\[
\Phi (0) - \mathcal {R} \Phi (x) \leqslant \varepsilon / 6 + 2 \Phi (0) \mathrm{Q} (x)\tag{63}
\]

for all \(x \in \mathbf{F}\) .

Set \(r = \left(12\Phi(0)\mathrm{Tr}(\mathbf{Q}/\mathbf{H})\varepsilon^{-1}\right)^{1/2}\) and denote by \(K\) the set of \(x' \in F_{\mathcal{T}}'\) such that \(\langle x, x' \rangle^2 \leqslant r^2\mathrm{H}(x)\) for all \(x \in F\) . Since \(H^{1/2}\) is a continuous semi-norm on \(F\) , the set \(K\) is equicontinuous and closed in \(F_{\mathcal{T}}'\) ; it is therefore compact in \(F_{\mathcal{T}}'\) by Ascoli's theorem (GT, X, §2, No.~5, Cor. 1 of Th. 2).

Let V be a closed linear subspace of \(F'_{T}\) with finite codimension; then, V is the orthogonal of a finite-dimensional linear subspace T of F. Let \(\mu_{V}\) be the measure on \(T'\) that is the image of the promeasure \(\mu\) on \(F'_{T}\) under the mapping \(p_{V}\) that is the transpose of the canonical injection of T into F; its Fourier transform is the restriction of \(\Phi\) to T. Finally, by the Hahn--Banach theorem (TVS, II, §3, No.~2, Cor. 1 of Th. 1), \(p_{\mathrm{V}}(\mathrm{K})\) is equal to the set \(C_{r}\) of \(x' \in T'\) such that \(\langle x, x' \rangle^{2} \leqslant r^{2}\mathrm{H}(x)\) for all \(x \in T\) . By the inequality (63), one can apply Prop. 10 of No.~9 to the measure \(\mu_{V}\) on \(T'\) , on taking for q the restriction of \(2\Phi(0)\mathrm{Q}\) to T and for h that of H. Then \(\mathrm{Tr}(q/h) \leqslant 2\Phi(0)\operatorname{Tr}(\mathrm{Q}/\mathrm{H})\) , whence

\[
\mu_ {\mathrm{V}} \left(\mathrm{T} ^ {\prime} - \mathrm{C} _ {r}\right) \leqslant 3 \left(\frac {\varepsilon}{6} + 2 \Phi (0) \operatorname{Tr} (\mathrm{Q} / \mathrm{H}) r ^ {- 2}\right) = \varepsilon .
\]

Since \(p_V\) defines, by passage to the quotient, an isomorphism of \(F_{\mathcal{T}}' / V\) onto \(T'\) , Prop. 1 of No.~1 then shows that \(\mu\) is a measure on \(F_{\mathcal{T}}'\) .

Q.E.D.

COROLLARY. --- Let F be a barreled nuclear space, T a locally convex topology on \(F'\) intermediate to \(T_{s}\) and \(T_{c}\) , \(\mu\) a promeasure on \(F'_{T}\) , and \(\Phi\) the Fourier transform of \(\mu\) . For \(\mu\) to be a measure, it is necessary and sufficient that \(\Phi\) be continuous on F.

Necessity follows from Prop. 9 of No.~8 and sufficiency from Th. 2.

Remark. --- Let F be a barreled space and T a locally convex topology on \(F'\) intermediate to \(T_s\) and \(T_c\) . Every subset of \(F'\) compact for T is compact for the coarser topology \(T_s\) . Conversely, let K be a subset of \(F'\) compact for \(T_s\) . Since F is barreled, K is equicontinuous (TVS, III, §4, No.~2, Th. 1); but by Ascoli's theorem, every equicontinuous subset of \(F'\) is relatively compact for \(T_c\) and a fortiori for T, therefore K is contained in a subset of \(F'\) compact for T. It is not difficult to infer from this that the identity mapping of \(F'_T\) onto \(F'_{T_s}\) defines a bijection between the sets of measures on these two spaces.

